% ============================================================
%  Conclusiones
% ============================================================
\section{Conclusiones}
\label{sec:conclusiones}

El objetivo del trabajo fue diseñar y construir un sistema semiautomatizado,
formado por un dispositivo y una interfaz de software, capaz de ejecutar en final
de línea los ensayos de seguridad eléctrica de la norma IEC~60335 con precisión,
repetibilidad y trazabilidad, y con menor intervención y riesgo para el operario.
Al cierre del trabajo, ese objetivo se cumple en forma completa para el ensayo de
consumo y en forma parcial para los ensayos de seguridad, cuya lógica está
terminada pero cuyo hardware de medición está pendiente.

\subsection{Logros respecto del objetivo y el alcance}

\begin{itemize}
    \item Se construyó un \textbf{sistema integrado y funcional}: el software
          detecta el hardware, comanda los relés y el contactor, lee los
          medidores, emite el veredicto, registra cada ensayo y lo exporta. El
          ensayo de consumo quedó validado de punta a punta sobre cargas reales,
          en sus modos de final de línea, perfil y cíclico, con una repetibilidad
          del \SI{0,22}{\percent}.
    \item Se \textbf{esquematizó el equipo completo de SYGlab}, que sirvió de
          referencia para la arquitectura y para los esquemas de conexión de los
          ensayos de seguridad. Sus registros (86 verificaciones con patrones
          entre 2023 y 2025) mostraron que la puesta a tierra era la medición más
          delicada, lo que orientó los criterios de validez del software nuevo.
    \item Los \textbf{criterios de aprobación quedaron alineados con la norma}.
          La revisión de la IEC~60335-1 mostró que los límites de partida no la
          reflejaban: la rigidez dieléctrica se evaluaba solo por tensión, la fuga
          pasiva tenía un límite fijo de \SI{1500}{\micro\ampere}, la puesta a
          tierra usaba un valor único y el consumo porcentajes fijos. Se
          adoptaron los ensayos de rutina del Anexo~A y las tablas de tolerancia
          de la cláusula~10, y los límites pasaron a depender de los atributos de
          cada producto.
    \item El \textbf{software pasó de un esqueleto a una aplicación completa}: la
          interfaz inicial tenía mayormente botones sin lógica y hoy integra las
          ocho pestañas previstas, con 158 pruebas automáticas y una integración
          continua que efectivamente verifica cada cambio.
    \item Se escribió el \textbf{firmware del microcontrolador}, que no existía:
          el módulo venía con un firmware de fábrica que no atendía a los
          medidores.
    \item Se priorizó la \textbf{seguridad del operario}: el sistema nunca emite
          un aprobado ante una lectura fallida, abre los relés ante cualquier
          error, y la parada de emergencia y el cierre de la aplicación detienen el
          ensayo en menos de un segundo.
\end{itemize}

\subsection{Dificultades encontradas y cómo se resolvieron}
\label{subsec:dificultades}

\paragraph{Límites normativos mal interpretados.}
Los valores de la tabla de cotas de la formulación inicial (\SI{125}{\milli\ohm}
para puesta a tierra, \SI{50}{\milli\ampere} para rigidez dieléctrica o
\SI{\pm 10}{\percent} fijo para consumo) no coincidían con la norma. En
particular, la corriente de la fuente de alta tensión se había tomado como límite
de aprobación, cuando el criterio de la norma es de \SI{5}{\milli\ampere}. Se resolvió con una lectura detallada de la
norma, distinguiendo los ensayos de tipo de los de rutina, y con pruebas
automáticas que fijan cada límite.

\paragraph{Puesta en marcha del microcontrolador.}
El ESP32 no era reconocido (el cable USB era solo de carga, y luego faltaba el
controlador del conversor CP2102), y una vez conectado al medidor entraba en un
bucle de reinicio. La causa era doble: el medidor estaba en los pines del puerto
USB y, alimentado con \SI{5}{\volt}, su salida de datos imponía \SI{5}{\volt} en
un pin del ESP32 y le impedía arrancar. Se resolvió moviendo el medidor al
segundo puerto serie y alimentándolo desde \SI{3,3}{\volt}. Los bornes de potencia
del medidor, sin serigrafía, se identificaron midiendo con multímetro.

\paragraph{Interferencia electromagnética en el control de los relés.}
La placa de relés USB se desconectaba de la PC al abrir el contactor con una carga
con motor; con la carga resistiva no ocurría. La placa implementa el USB por
software en un microcontrolador pequeño, muy sensible a las perturbaciones que
genera la apertura de una carga inductiva. Se mitigó en dos frentes: el software
pasó a tolerar la desconexión (reintenta, confirma que los relés quedaron abiertos
o pide cortar la red) y se aprovechó que, al reiniciarse, la placa deja los relés
abiertos, lo que hace que la falla sea segura; además, se agregó un varistor en la
bobina del contactor, que redujo el problema sin eliminarlo. Fue la principal
lección de diseño del trabajo: los supresores en las bobinas, la aislación
galvánica del USB y la elección de controladores con USB por hardware deben
considerarse desde el inicio.

\paragraph{Subestimación del consumo.}
La prueba de punta a punta reveló que el consumo se subestimaba en un
\SI{27}{\percent} porque el promedio incluía las lecturas del transitorio de cierre
del contactor, lo que hacía fallar a un aparato sano. Se corrigió descartando las
lecturas de estabilización y se verificó en el banco. Esta dificultad muestra el
valor de probar con el hardware real: las pruebas con hardware simulado no podían
detectarla.

\paragraph{Calidad del proceso de software.}
La integración continua estaba mal configurada desde el primer commit y nunca
había ejecutado las pruebas ni el análisis estático; además, las correcciones de
calibración configuradas nunca se aplicaban y la API tenía credenciales de fábrica
en texto plano. Se corrigieron los tres puntos: la integración continua quedó en
verde, las correcciones se aplican a todas las lecturas y las contraseñas se
guardan con hash, rechazando las de fábrica.

\paragraph{Plazos de los componentes.}
Las demoras en los envíos de algunos componentes retrasaron unos días la
construcción del prototipo. El impacto se compensó con el adelanto del software:
el uso de herramientas de programación asistida por inteligencia artificial
permitió completar la implementación y las pruebas alrededor de dos semanas antes
de lo previsto, y ese tiempo se reinvirtió en la integración y en las pruebas en
el banco (Anexo~\ref{anx:cumplimiento_cronograma}).

\subsection{Limitaciones}

\begin{itemize}
    \item Los ensayos de fuga, puesta a tierra y rigidez dieléctrica del software
          nuevo no se ejecutaron con hardware real: el prototipo propio no incluye
          su hardware de medición, y en el equipo de SYGlab solo hay registros del
          sistema anterior.
    \item No se contó con un instrumento patrón: los resultados no constituyen
          una validación metrológica y las correcciones de calibración no están
          ajustadas.
    \item El banco mide a la tensión de red disponible y no a la tensión nominal
          del aparato, lo que afecta el veredicto de aparatos cercanos al límite.
    \item La interferencia de la placa de relés está mitigada, pero no eliminada.
    \item El banco no tiene interruptor diferencial.
\end{itemize}

\subsection{Trabajo futuro}

\begin{itemize}
    \item Completar el hardware de los ensayos de seguridad: medidor de
          \SI{10}{\ampere} con red de medición IEC~60990 para fuga, transformador
          de puesta a tierra verificado a \SI{10}{\ampere} y transformador de alta
          tensión con medición de \SI{5}{\milli\ampere}.
    \item Calibrar el sistema contra un instrumento patrón y validar con los
          patrones del modo de verificación.
    \item Comandar el contactor desde el ESP32 con un relé optoacoplado (el
          firmware ya admite los comandos) y aislar galvánicamente el USB.
    \item Incorporar un interruptor diferencial y regular o corregir la tensión de
          ensayo de consumo.
    \item Registrar en el historial la tensión y la corriente medidas en rigidez
          dieléctrica, y medir los tiempos frente al método manual para
          confirmar las estimaciones de tiempo e intervención del operario.
\end{itemize}

\subsection{Reflexión final}

Más allá del equipo construido, el trabajo dejó un aprendizaje sobre el proceso
de desarrollo de un instrumento de seguridad: leer la norma en detalle antes de
fijar un límite, desconfiar de un resultado hasta comprobarlo con el hardware real
y diseñar cada parte para que, ante una falla, el sistema quede en un estado
seguro. El prototipo ofrece una base abierta, de bajo costo y verificada para
completar los ensayos restantes y avanzar hacia su uso en una línea de producción.

En lo personal, este trabajo me permitió recorrer el ciclo completo de un
desarrollo de ingeniería: desde la formulación y la lectura de la norma hasta el
armado del banco, la programación y la verificación con cargas reales. Las
etapas que más me enseñaron fueron las que no salieron como estaba previsto: el
microcontrolador que no arrancaba, la placa de relés que se desconectaba o el
ensayo que rechazaba un aparato sano. Cada una obligó a volver a las mediciones,
plantear hipótesis y comprobarlas antes de dar algo por resuelto, y me mostró que
en un equipo de seguridad la confiabilidad se construye tanto en el diseño como
en la verificación. También aprendí a integrar herramientas de programación
asistida por inteligencia artificial sin delegar el criterio técnico: aceleraron
el desarrollo, pero cada límite normativo y cada resultado se contrastaron con la
norma y con el banco. Cierro este trabajo con un prototipo funcional, con una
idea clara de lo que falta para completarlo y con la convicción de que la
automatización de los ensayos de seguridad es un camino viable y necesario para
la industria local.
